\section{Specificaties}




\begin{table}[htbp]
    \centering
    \caption{Specificaties en Eisen}
    \label{tab:specificaties}
    \begin{tabular}{|p{4cm}|p{10cm}|}
        \hline
        \textbf{Specificatie} & \textbf{Onderbouwing / Eis} \\
        \hline \hline
        Snelheid & Het systeem moet snelheden tot minstens 80 km/u accuraat registreren. Dit dekt de topsnelheid in de paardensport inclusief een veiligheidsmarge van 10\%. \\
        \hline
        Nauwkeurigheid (Open Veld) & De toevoeging van sensor fusion mag de baseline RTK-precisie (typisch $<5$ cm) niet negatief beïnvloeden. \\
        \hline
        Dimensionaliteit & Kijkend naar de projectafbakening beperkt de analyse en positiebepaling zich strikt tot het 2D vlak (Horizontale $X, Y$ positie en Heading). \\
        \hline
        Foutdetectie & Het systeem moet 'False Fixed' oplossingen detecteren. Een afwijking tussen de GNSS-meting en de IMU-voorspelling van meer dan 15 cm (binnen één tijdstap) moet gemarkeerd worden als onbetrouwbaar. Idealiter in de enkele centimeters. \\
        \hline
        Update Frequentie & Het huidige systeemstandaard bedraagt 500 ms waarbij het gewenst is van het systeem dat de output minimaal binnen deze tijd moet plaatsvinden. \\
        \hline
        Verwerkingssnelheid & De dataverwerking moet real-time plaatsvinden. De berekeningstijd per stap moet korter zijn dan het sample interval om vertraging te voorkomen. \\
        \hline
        Communicatie Protocol & Output zal gelijk moeten zijn aan het X2 Link protocol van MYLAPS. \\
        \hline
    \end{tabular}
\end{table}

\newpage